% ==================== THEME 1 — VARIABLES SIMPLES ====================

% ---- var-declaration ----
\element{bank-var-declaration}{
    \begin{question}{var-declaration-v1}
        Quelle ligne déclare correctement une variable entière nommée \texttt{age} ?
        \begin{multicols}{4}
            \begin{reponses}
                \bonne{\lstinline[language=C]|int age;|}
                \mauvaise{\lstinline[language=C]|int age|}
                \mauvaise{\lstinline[language=C]|age int;|}
                \mauvaise{\lstinline[language=C]|integer age;|}
            \end{reponses}
        \end{multicols}
    \end{question}
}
\element{bank-var-declaration}{
    \begin{question}{var-declaration-v2}
        Quelle ligne déclare correctement une variable entière nommée \texttt{score} ?
        \begin{multicols}{4}
            \begin{reponses}
                \bonne{\lstinline[language=C]|int score;|}
                \mauvaise{\lstinline[language=C]|int score|}
                \mauvaise{\lstinline[language=C]|score int;|}
                \mauvaise{\lstinline[language=C]|integer score;|}
            \end{reponses}
        \end{multicols}
    \end{question}
}
\element{bank-var-declaration}{
    \begin{question}{var-declaration-v3}
        Quelle ligne déclare correctement une variable entière nommée \texttt{compteur} ?
        \begin{multicols}{4}
            \begin{reponses}
                \bonne{\lstinline[language=C]|int compteur;|}
                \mauvaise{\lstinline[language=C]|int compteur|}
                \mauvaise{\lstinline[language=C]|compteur int;|}
                \mauvaise{\lstinline[language=C]|integer compteur;|}
            \end{reponses}
        \end{multicols}
    \end{question}
}

\element{variables}{
    \restituegroupe[1]{bank-var-declaration}
}

% ---- var-type-decimal ----
\element{bank-var-type-decimal}{
    \begin{question}{var-type-decimal-v1}
        Pour stocker un nombre à virgule comme \texttt{3.14}, quel type utiliser en priorité ?
        \begin{multicols}{4}
            \begin{reponses}
                \mauvaise{\lstinline[language=C]|int|}
                \bonne{\lstinline[language=C]|float|}
                \mauvaise{\lstinline[language=C]|char|}
                \mauvaise{\lstinline[language=C]|string|}
            \end{reponses}
        \end{multicols}
    \end{question}
}
\element{bank-var-type-decimal}{
    \begin{question}{var-type-decimal-v2}
        Pour stocker un nombre à virgule comme \texttt{2.5}, quel type utiliser en priorité ?
        \begin{multicols}{4}
            \begin{reponses}
                \mauvaise{\lstinline[language=C]|int|}
                \bonne{\lstinline[language=C]|float|}
                \mauvaise{\lstinline[language=C]|char|}
                \mauvaise{\lstinline[language=C]|string|}
            \end{reponses}
        \end{multicols}
    \end{question}
}
\element{bank-var-type-decimal}{
    \begin{question}{var-type-decimal-v3}
        Pour stocker un nombre à virgule comme \texttt{9.99}, quel type utiliser en priorité ?
        \begin{multicols}{4}
            \begin{reponses}
                \mauvaise{\lstinline[language=C]|int|}
                \bonne{\lstinline[language=C]|float|}
                \mauvaise{\lstinline[language=C]|char|}
                \mauvaise{\lstinline[language=C]|string|}
            \end{reponses}
        \end{multicols}
    \end{question}
}

\element{variables}{
    \restituegroupe[1]{bank-var-type-decimal}
}

% ---- var-type-caractere ----
\element{bank-var-type-caractere}{
    \begin{question}{var-type-caractere-v1}
        Quel type permet de stocker un unique caractère, comme \texttt{'A'} ?
        \begin{multicols}{4}
            \begin{reponses}
                \mauvaise{\lstinline[language=C]|int|}
                \mauvaise{\lstinline[language=C]|float|}
                \bonne{\lstinline[language=C]|char|}
                \mauvaise{\lstinline[language=C]|void|}
            \end{reponses}
        \end{multicols}
    \end{question}
}
\element{bank-var-type-caractere}{
    \begin{question}{var-type-caractere-v2}
        Quel type permet de stocker un unique caractère, comme \texttt{'Z'} ?
        \begin{multicols}{4}
            \begin{reponses}
                \mauvaise{\lstinline[language=C]|int|}
                \mauvaise{\lstinline[language=C]|float|}
                \bonne{\lstinline[language=C]|char|}
                \mauvaise{\lstinline[language=C]|void|}
            \end{reponses}
        \end{multicols}
    \end{question}
}
\element{bank-var-type-caractere}{
    \begin{question}{var-type-caractere-v3}
        Quel type permet de stocker un unique caractère, comme \texttt{'x'} ?
        \begin{multicols}{4}
            \begin{reponses}
                \mauvaise{\lstinline[language=C]|int|}
                \mauvaise{\lstinline[language=C]|float|}
                \bonne{\lstinline[language=C]|char|}
                \mauvaise{\lstinline[language=C]|void|}
            \end{reponses}
        \end{multicols}
    \end{question}
}

\element{variables}{
    \restituegroupe[1]{bank-var-type-caractere}
}

% ---- var-nommage ----
% Durci : le piège "commence par un chiffre" est repéré instantanément par tous
% les débutants ; on le remplace par le piège plus subtil du mot-clé réservé
% utilisé comme nom de variable, et on ajoute un identifiant à underscore
% initial (valide, mais souvent cru interdit) pour forcer une vraie analyse.
\element{bank-var-nommage}{
    \begin{question}{var-nommage-v1}
        Quel nom de variable est \textbf{invalide} en C ?
        \begin{multicols}{5}
            \begin{reponses}
                \mauvaise{\lstinline[language=C]|nb_essais|}
                \mauvaise{\lstinline[language=C]|_compteur|}
                \bonne{\lstinline[language=C]|for|}
                \mauvaise{\lstinline[language=C]|essai2|}
                \mauvaise{\lstinline[language=C]|Essai_Final|}
            \end{reponses}
        \end{multicols}
    \end{question}
}
\element{bank-var-nommage}{
    \begin{question}{var-nommage-v2}
        Quel nom de variable est \textbf{invalide} en C ?
        \begin{multicols}{5}
            \begin{reponses}
                \mauvaise{\lstinline[language=C]|nb_tours|}
                \mauvaise{\lstinline[language=C]|_limite|}
                \bonne{\lstinline[language=C]|while|}
                \mauvaise{\lstinline[language=C]|tour3|}
                \mauvaise{\lstinline[language=C]|TourMax|}
            \end{reponses}
        \end{multicols}
    \end{question}
}
\element{bank-var-nommage}{
    \begin{question}{var-nommage-v3}
        Quel nom de variable est \textbf{invalide} en C ?
        \begin{multicols}{5}
            \begin{reponses}
                \mauvaise{\lstinline[language=C]|note_max|}
                \mauvaise{\lstinline[language=C]|_seuil|}
                \bonne{\lstinline[language=C]|if|}
                \mauvaise{\lstinline[language=C]|note1|}
                \mauvaise{\lstinline[language=C]|NoteFinale|}
            \end{reponses}
        \end{multicols}
    \end{question}
}

\element{variables}{
    \restituegroupe[1]{bank-var-nommage}
}

% ---- var-initialisation ----
\element{bank-var-initialisation}{
    \begin{question}{var-initialisation-v1}
        Que vaut \texttt{x} juste après l'exécution de la ligne \lstinline[language=C]|int x;| (sans affectation) ?
        \begin{multicols}{4}
            \begin{reponses}
                \mauvaise{0}
                \bonne{Une valeur indéterminée}
                \mauvaise{-1}
                \mauvaise{Une erreur de compilation}
            \end{reponses}
        \end{multicols}
    \end{question}
}
\element{bank-var-initialisation}{
    \begin{question}{var-initialisation-v2}
        Que vaut \texttt{n} juste après l'exécution de la ligne \lstinline[language=C]|int n;| (sans affectation) ?
        \begin{multicols}{4}
            \begin{reponses}
                \mauvaise{0}
                \bonne{Une valeur indéterminée}
                \mauvaise{-1}
                \mauvaise{Une erreur de compilation}
            \end{reponses}
        \end{multicols}
    \end{question}
}
\element{bank-var-initialisation}{
    \begin{question}{var-initialisation-v3}
        Que vaut \texttt{total} juste après l'exécution de la ligne \lstinline[language=C]|int total;| (sans affectation) ?
        \begin{multicols}{4}
            \begin{reponses}
                \mauvaise{0}
                \bonne{Une valeur indéterminée}
                \mauvaise{-1}
                \mauvaise{Une erreur de compilation}
            \end{reponses}
        \end{multicols}
    \end{question}
}

\element{variables}{
    \restituegroupe[1]{bank-var-initialisation}
}
